% §16 Problem 1: velocities of the two particles in the L system, obtained by
% adding V to their opposite velocities in the C system (cf. Fig. 14 of the book)
\documentclass{SolutionFigure}

\begin{document}
\begin{tikzpicture}[scale=2.2]
    % Parameters
    \def\V{1.3}     % velocity V of the primary particle
    \def\va{1}      % speed v_10 of particle 1 in the C system
    \def\vb{1.6}    % speed v_20 of particle 2 in the C system
    \def\th{55}     % angle theta_0

    \coordinate (A) at (0,0);
    \coordinate (T) at (\V,0);
    \coordinate (X) at ({\V+0.9},0);
    \coordinate (P1) at ($(T)+(\th:\va)$);
    \coordinate (P2) at ($(T)+({\th+180}:\vb)$);

    % Direction of V and of the emission in the C system
    \draw[guide] (A) -- (X);
    \draw[guide] (P2) -- ($(T)!1.25!(P1)$);
    \node[annotation, anchor=west] at (X) {$x$};

    % Velocities in the C system, from the tip of V
    \draw[vector] (A) -- node[below=1pt, pos=0.8] {$\mathbf{V}$} (T);
    \draw[vector] (T) -- node[right=1pt, pos=0.6] {$\mathbf{v}_{10}$} (P1);
    \draw[vector] (T) -- node[right=2pt, pos=0.55] {$\mathbf{v}_{20}$} (P2);

    % Velocities in the L system
    \draw[vector, highlight] (A) -- node[above left=0pt, pos=0.6] {$\mathbf{v}_1$} (P1);
    \draw[vector, highlight] (A) -- node[left=1pt, pos=0.6] {$\mathbf{v}_2$} (P2);

    % Angles: theta_0 = theta_10, and theta_20 = pi - theta_0 between V and v_20
    \pic[angle arc, angle radius=5mm, angle eccentricity=1.45, "$\theta_0$"] {angle = X--T--P1};
    \pic[angle arc, angle radius=4mm, angle eccentricity=1.9, "$\theta_{20}$"] {angle = P2--T--X};
    \pic[angle arc, angle radius=9mm, angle eccentricity=1.3, "$\theta_1$"] {angle = T--A--P1};
    \pic[angle arc, angle radius=6mm, angle eccentricity=1.4, "$\theta_2$"] {angle = P2--A--T};
\end{tikzpicture}
\end{document}
